\addtocontents{toc}{\protect\Needspace{3\baselineskip}}
\originalchapter{第 26 讲公开期末复习题}
三题均实际出现在 Bryan Owens 整理的 MIT 8.251 第 26 讲笔记. 前两题附学生笔记解法, 本稿重新推导并改正个别省略与排印; 第三题的公开文件没有给出解答, 以下独立补写. 这三题是覆盖冻结后的新增发现, 在补充清单中保留其来源与编号, 不混入正式期末考试题数.

\src{8.251 Lecture26 Review1, 原 PDF 第 1--2 页.}
\begin{exercise}
标准充满空间的玻色开弦, 横向 Virasoro 模为 $L_n^\perp=\frac12\sum_k:\alpha_{n-k}^J\alpha_k^J:$, 零横向动量真空满足 $\alpha_0|0\rangle=0$.
(a) 用正规序振子写 $L_{-6}^\perp|0\rangle$.
(b) 求 $L_{-6}^\perp\alpha_{-1}^I|0\rangle$.
\end{exercise}
\begin{solution}
只有 $k=-1,-2,-3,-4,-5$ 的双创建项能留下. 交换相同模式并合并重复配对得
\[
 L_{-6}|0\rangle=
 \left(\alpha_{-1}\cdot\alpha_{-5}
 +\alpha_{-2}\cdot\alpha_{-4}
 +\tfrac12\alpha_{-3}\cdot\alpha_{-3}\right)|0\rangle.
\]
零模项为零, 含正频湮灭算符的正规序项杀真空. 由 $[L_m,\alpha_n^I]=-n\alpha_{m+n}^I$,
\[
 L_{-6}\alpha_{-1}^I|0\rangle
 =\alpha_{-1}^I\left(\alpha_{-1}\cdot\alpha_{-5}
 +\alpha_{-2}\cdot\alpha_{-4}
 +\tfrac12\alpha_{-3}^2\right)|0\rangle
 +\alpha_{-7}^I|0\rangle.
\]
原笔记第二式的一处简写丢掉了第一个项中的 $\alpha_{-1}^I$, 本稿保留它. 两边总层数均为 7, 是一个立即可用的检查.
\end{solution}

\src{8.251 Lecture26 Review2, 原 PDF 第 1、3--4 页.}
\begin{exercise}
在 $D=26$ 时空, D2 膜位于 $z=0$, 沿 $x,y$; D1 膜沿 $x$, 位于 $y=0,z=\pi\sqrt{\alpha'}$. 其他空间坐标均同为零. 考虑从 D2 到 D1 的开玻色弦.
(a) 用横向振子层数 $N$ 写质量式.
(b) 描述基态与最低质量层.
(c) 求 $f(q)=\sum_r d(r)q^r$, 其中 $r=\alpha'M^2$.
\end{exercise}
\begin{solution}
将光锥取在共同的 $t,x$ 方向, $y$ 是一个 ND 坐标, $z$ 是伸长的 DD 坐标, 其余 22 个横向方向也是 DD. 因而有 23 个整数频率与一个半整数频率, 没有连续横向零模动量. 两端距离 $b=\pi\sqrt{\alpha'}$ 给
$\alpha'(T_sb)^2=1/4$, 零点能为 $-1+1/16=-15/16$. 所以
\[
 \alpha'M^2=N-\frac{11}{16}.
\]
基态层 $N=0$ 有一态, 质量平方 $-11/(16\alpha')$. $N=1/2$ 是唯一半整数第一振子的一个态, 仍为快子, $\alpha'M^2=-3/16$. $N=1$ 有 23 个整数第一模式, 另有半整数第一振子激发两次的一态, 共 24 态, $\alpha'M^2=5/16$. 半整数振子在本题是玻色, 不能用 Pauli 排除禁止重复激发.

更高层独立模式相乘,
\[
 f(q)=q^{-11/16}\prod_{n=1}^\infty
 \frac1{(1-q^n)^{23}(1-q^{n-1/2})}.
\]
低层系数以 $N$ 为指数为 $1,1,24,25,324$（层 $0,1/2,1,3/2,2$）, 随包有限截断枚举核验. 态沿共同 $x$ 方向传播, 用 $(p^+,p^-)$ 描述其 $1+1$ 维运动, 不误说它有 $y,z$ 的自由动量.
\end{solution}

\src{8.251 Lecture26 Review3, 原 PDF 第 1--2 页; 无公开解答.}
\begin{exercise}
形式上考虑 $D=18$ 的 R 扇区超弦横向振子.
(a) 给横向指标范围, 将包含玻色与整数频费米振子的朴素对称序质量式改为精确正规序形式.
(b) 写未作 GSO 投影的生成函数 $f_R(q)$.
\end{exercise}
\begin{solution}
横向指标 $I=2,\ldots,17$, 共 16 个. 每方向玻色有限零点能 $-1/24$ 与 R 费米的 $+1/24$ 抵消, 零模 Clifford 振子不增加层数. 所以形式质量式为
\[
 \alpha'M^2=\sum_{n>0,I}\alpha_{-n}^I\alpha_n^I
 +\sum_{n>0,I}n d_{-n}^Id_n^I.
\]
16 个实 R 零模可组合成八对复费米零模, 未投影真空空间维数 $2^8=256$. 非零频率每层有 16 个可重复的玻色模式与 16 个只能占据一次的费米模式, 故
\[
 f_R(q)=256\prod_{n=1}^\infty
 \left(\frac{1+q^n}{1-q^n}\right)^{16}.
\]
若单独作手征投影, 真空数会减半, 但原题没有要求它, 因而不擅自把 256 改为 128. 这只是规定振子代数下的形式谱计数, 标准平直 RNS 超弦临界维数是 10, $D=18$ 未满足全体 Lorentz 与 Weyl 一致性; 本题答案不表示存在已构造的一致十八维超弦理论.
\end{solution}
